\documentclass[12pt,a4paper]{report}
% Общая преамбула. Подключается из subject/main.tex после \documentclass.

% ============================================================
% Базовые шрифты и кодировка
% ============================================================
\usepackage[T2A]{fontenc}
\usepackage[utf8]{inputenc}
\usepackage[russian]{babel}

\renewcommand{\rmdefault}{cmr}
\renewcommand{\sfdefault}{cmss}
\renewcommand{\ttdefault}{cmtt}
\renewcommand{\familydefault}{\rmdefault}

% ============================================================
% Математика и базовое оформление
% ============================================================
\usepackage{amsmath,amssymb,amsthm,mathtools}
\usepackage{mathrsfs}
\usepackage{bbold}
\usepackage{enumitem}
\usepackage{xparse}
\usepackage{xcolor}
\usepackage{microtype}
\usepackage{geometry}          % для настройки полей
\usepackage{parskip}           % для расстояний между абзацами (без параметров)
\usepackage{hyperref}
\usepackage[nameinlink,noabbrev]{cleveref}
\usepackage{tikz}
\usetikzlibrary{
    positioning,
    arrows.meta,
    calc,
    intersections,
    patterns,
    shapes.geometric,
    decorations.markings,
    matrix,
    fit,
    backgrounds
}

% ============================================================
% Рисунки, графики, схемы, таблицы и фрагменты кода
% ============================================================
\usepackage{graphicx}      % изображения
\usepackage{float}         % точное размещение рисунков/таблиц через [H]
\usepackage{caption}       % подписи к рисункам и таблицам
\usepackage{subcaption}    % несколько рисунков внутри одной фигуры
\usepackage{wrapfig}       % обтекание рисунков текстом
\usepackage{pgfplots}      % графики функций и численных данных
\pgfplotsset{compat=1.18}
\usepackage{tikz-cd}       % коммутативные диаграммы

\usepackage{booktabs}      % аккуратные таблицы
\usepackage{array}
\usepackage{tabularx}
\usepackage{longtable}
\usepackage{multirow}

\usepackage{listings}      % код: C/C++/Python/SQL и др.

% Дополнительная математика
\usepackage{bm}            % жирные математические символы: \bm{x}
\usepackage{bbm}           % красивый blackboard bold для цифр: \mathbbm{0}, \mathbbm{1}
\usepackage{esint}         % \oiint, \varoiint и другие интегралы
\usepackage{cancel}        % \cancel{}, \bcancel{}, \xcancel{}

% ============================================================
% Отступы и расстояния
% ============================================================
\geometry{a4paper,top=2cm,bottom=2cm,left=3cm,right=3cm,marginparwidth=1.75cm}
% parskip уже загружен без параметров, что даёт нулевой отступ первой строки
% и положительный межстрочный интервал между абзацами.

\setlength{\emergencystretch}{3em}

\hypersetup{
    colorlinks=true,
    linkcolor=black,
    urlcolor=black,
    citecolor=black,
    linktoc=all
}

% ============================================================
% Доказательства
% ============================================================
\addto\captionsrussian{%
    \renewcommand{\proofname}{Доказательство}%
}

\makeatletter
\renewenvironment{proof}[1][\proofname]{%
    \par
    \begingroup
    \leftskip=\dimexpr-\@totalleftmargin\relax
    \rightskip=0pt
    \parindent=0pt
    \noindent\textit{\underline{#1.}}\par\nobreak
    \noindent\ignorespaces
}{%
    \unskip\nobreak\hspace{0.15em}\ensuremath{\blacksquare}%
    \par
    \endgroup
}
\makeatother

% ============================================================
% Заголовки
% ============================================================
\usepackage{titlesec}
\titleclass{\chapter}{straight}

\titleformat{\chapter}[display]
  {\normalfont\huge\bfseries\raggedright}
  {\chaptername\ \thechapter}
  {0.7em}
  {\Huge\raggedright}

\titleformat{\section}
  {\normalfont\Large\bfseries}
  {\thesection}{0.7em}{}

\titleformat{\subsection}
  {\normalfont\large\bfseries}
  {\thesubsection}{0.7em}{}

\titleformat{\subsubsection}
  {\normalfont\normalsize\bfseries}
  {\thesubsubsection}{0.7em}{}

\titlespacing*{\chapter}{0pt}{0.45em}{0.45em}
\titlespacing*{\section}{0pt}{0.55em}{0.25em}
\titlespacing*{\subsection}{0pt}{0.45em}{0.18em}
\titlespacing*{\subsubsection}{0pt}{0.35em}{0.15em}

% ============================================================
% Лекции
% ============================================================
\newcounter{lection}
\NewDocumentCommand{\newlection}{m}{%
    \refstepcounter{lection}%
    \par\smallskip
    {\centering\color{gray!70}\normalsize Лекция \thelection. #1\par}
    \vspace{0.1em}
}

% ============================================================
% Теоремы, определения, замечания
% ============================================================
\makeatletter
\NewDocumentCommand{\definition}{o +m}{%
    \par
    \noindent
    \ifdim\@totalleftmargin>0pt
        \hspace*{\dimexpr-\@totalleftmargin\relax}%
    \fi
    \textbf{Определение\IfValueT{#1}{ (#1)}.}~#2
    \par
}
\makeatother

\ExplSyntaxOn

\NewDocumentCommand{\theorem}{o +m}{%
    \par\noindent
    \textbf{%
        Теорема
        \IfValueT{#1}{%
            \regex_match:nnTF { \A \d+ \Z } { #1 }
                {~#1}
                {~(#1)}
        }%
        .
    }~#2
    \par
}

\ExplSyntaxOff

\NewDocumentCommand{\proposal}{o +m}{%
    \par\noindent
    \textbf{Предложение\IfValueT{#1}{~#1}.}~#2
    \par
}

\NewDocumentCommand{\fact}{o +m}{%
    \par\noindent
    \textit{Факт\IfValueT{#1}{ (#1)}.}~#2
    \par
}

\NewDocumentCommand{\question}{+m}{%
    \par\noindent
    \textbf{Вопрос.}~#1
    \par
}

\NewDocumentCommand{\note}{+m}{%
    \par\noindent
    \textbf{Замечание.}~#1
    \par
}

\NewDocumentCommand{\corollary}{o +m}{%
    \par\noindent
    \textbf{Следствие\IfValueT{#1}{~#1}.}~#2
    \par
}

\NewDocumentCommand{\example}{+m}{%
    \par\noindent
    \textbf{Пример.}~#1
    \par
}

% ============================================================
% Списки
% ============================================================
\NewDocumentCommand{\numbers}{+m}{%
    \begin{enumerate}[leftmargin=2em,itemsep=0.25em,topsep=0.3em]
    #1
    \end{enumerate}
}

\NewDocumentCommand{\bullets}{+m}{%
    \begin{itemize}[leftmargin=2em,itemsep=0.25em,topsep=0.3em]
    #1
    \end{itemize}
}

\NewDocumentCommand{\provehere}{+m}{%
    \par\smallskip
    \begin{proof}#1\end{proof}
}

\NewDocumentCommand{\provewthen}{m +m}{%
    \par\smallskip
    \begin{proof}[#1]#2\end{proof}
}

\NewDocumentCommand{\provebullets}{+m}{%
    \begin{proof}
    \begin{itemize}[leftmargin=2em,itemsep=0.25em,topsep=0.3em]
    #1
    \end{itemize}
    \end{proof}
}

\NewDocumentCommand{\provenumbers}{+m}{%
    \begin{proof}
    \begin{enumerate}[leftmargin=2em,itemsep=0.25em,topsep=0.3em]
    #1
    \end{enumerate}
    \end{proof}
}

% ============================================================
% Дополнительные команды
% ============================================================
\usepackage{empheq}      % для \encircle
\usepackage{changepage}  % для \blockindent

% Леммы
\newtheorem{lemma_env}{Лемма}[section]
\newcommand{\lemma}[2][]{\begin{lemma_env}[#1]#2\end{lemma_env}}

% Отступ (переименовано из \indent, чтобы не ломать ядро LaTeX)
\newcommand{\blockindent}[1]{%
  \begin{adjustwidth}{1.5cm}{}#1\end{adjustwidth}%
}

% Одностраничные блоки
\newcommand{\singlepage}[1]{\begin{samepage}#1\end{samepage}}

% Сборки формул
\renewcommand{\gather}[1]{\begin{gather*}#1\end{gather*}}
\renewcommand{\multline}[1]{\begin{multline*}#1\end{multline*}}

% Обрамлённые формулы
\newcommand{\encircle}[1]{\begin{empheq}[box=\fbox]{align*}#1\end{empheq}}

% Математические операторы и символы
\newcommand{\proddots}{\cdot\ldots\cdot}
\newcommand{\abs}[1]{\left|#1\right|}
\newcommand{\angles}[1]{\left\langle#1\right\rangle}
\newcommand{\der}[2]{\frac{\partial #1}{\partial #2}}
\DeclareMathOperator{\differential}{d}
\newcommand{\dif}{\differential\!}
\newcommand{\fmax}{\lor}
\newcommand{\fmin}{\land}
\DeclareMathOperator*{\esssup}{ess\,sup}
\DeclareMathOperator*{\essssup}{ess\,sup}
\DeclareMathOperator*{\esssssup}{ess\,sup}
\newcommand{\Map}{\longrightarrow}
\newcommand{\almost}[1]{\underset{#1}{\overset{\textup{почти всюду}}\Map}}
\newcommand{\circlesign}[1]{\mathrel{\tikz[baseline=-0.5ex]{\node[draw, circle, inner sep=1pt] (char) {\vphantom{$-$}$#1$};}}}

% Дополнительные команды для доказательств
\newcommand{\indentlemma}[3][]{\blockindent{\lemma[#1]{#2\prove[Доказательство леммы]{#3}}}}
\newcommand{\prove}[2][Доказательство]{\begin{proof}[#1]#2\end{proof}}

% ============================================================
% Часто используемые математические обозначения
% ============================================================
\newcommand{\N}{\mathbb{N}}
\newcommand{\Z}{\mathbb{Z}}
\newcommand{\Q}{\mathbb{Q}}
\newcommand{\R}{\mathbb{R}}
\renewcommand{\C}{\mathbb{C}}
\newcommand{\Ord}{\mathrm{Ord}}
\newcommand{\Card}{\mathrm{Card}}

\renewcommand{\o}{\varnothing}
\newcommand{\bs}{\setminus}
\newcommand{\inv}{^{\complement}}
\newcommand{\then}{\Rightarrow}
\newcommand{\map}{\to}
\newcommand{\bydef}{\coloneqq}
\newcommand{\all}[1]{\left\{\begin{aligned}#1\end{aligned}\right.}
\newcommand{\any}[1]{\left[\begin{aligned}#1\end{aligned}\right.}
\newcommand{\exist}{\exists}
\newcommand{\dom}{\operatorname{dom}}
\newcommand{\rng}{\operatorname{rng}}
\newcommand{\defset}[2]{\left\{#1\;\middle|\;#2\right\}}
\newcommand{\switch}[1]{\left\{\begin{aligned}#1\end{aligned}\right.}

% ------------------------------------------------------------
% Большие операторы: индексы автоматически ставятся сверху/снизу
% ------------------------------------------------------------
\let\oldlim\lim
\renewcommand{\lim}{\oldlim\limits}
\let\oldsup\sup
\renewcommand{\sup}{\oldsup\limits}
\let\oldinf\inf
\renewcommand{\inf}{\oldinf\limits}
\let\oldmax\max
\renewcommand{\max}{\oldmax\limits}
\let\oldmin\min
\renewcommand{\min}{\oldmin\limits}

\let\oldsum\sum
\renewcommand{\sum}{\oldsum\limits}
\let\oldprod\prod
\renewcommand{\prod}{\oldprod\limits}
\let\oldcoprod\coprod
\renewcommand{\coprod}{\oldcoprod\limits}

\let\oldbigcap\bigcap
\renewcommand{\bigcap}{\oldbigcap\limits}
\let\oldbigcup\bigcup
\renewcommand{\bigcup}{\oldbigcup\limits}
\let\oldbigsqcup\bigsqcup
\renewcommand{\bigsqcup}{\oldbigsqcup\limits}

\let\oldbigvee\bigvee
\renewcommand{\bigvee}{\oldbigvee\limits}
\let\oldbigwedge\bigwedge
\renewcommand{\bigwedge}{\oldbigwedge\limits}
\let\oldbigoplus\bigoplus
\renewcommand{\bigoplus}{\oldbigoplus\limits}
\let\oldbigotimes\bigotimes
\renewcommand{\bigotimes}{\oldbigotimes\limits}
\let\oldbigodot\bigodot
\renewcommand{\bigodot}{\oldbigodot\limits}
\let\oldbiguplus\biguplus
\renewcommand{\biguplus}{\oldbiguplus\limits}

% Красивые blackboard-bold цифры (\mathbb не во всех шрифтах содержит цифры)
\newcommand{\bbzero}{\mathbbm{0}}
\newcommand{\bbone}{\mathbbm{1}}

% Универсальные обозначения для разных курсов
\newcommand{\norm}[1]{\left\lVert #1 \right\rVert}
\newcommand{\floor}[1]{\left\lfloor #1 \right\rfloor}
\newcommand{\ceil}[1]{\left\lceil #1 \right\rceil}
\newcommand{\paren}[1]{\left( #1 \right)}
\newcommand{\brackets}[1]{\left[ #1 \right]}
\newcommand{\set}[1]{\left\{ #1 \right\}}
\newcommand{\conj}[1]{\overline{#1}}
\newcommand{\vect}[1]{\bm{#1}}

\DeclareMathOperator{\RePart}{Re}
\DeclareMathOperator{\ImPart}{Im}
\DeclareMathOperator{\Arg}{Arg}
\DeclareMathOperator{\Ker}{Ker}
\DeclareMathOperator{\Image}{Im}
\DeclareMathOperator{\rank}{rank}
\DeclareMathOperator{\tr}{tr}
\DeclareMathOperator{\sgn}{sgn}
\DeclareMathOperator{\lcm}{lcm}
\DeclareMathOperator{\Span}{span}

% Differential notation in math; preserve the text dot-below accent.
\let\textdotbelow\d
\DeclareRobustCommand{\d}{\ifmmode\mathrm{d}\else\expandafter\textdotbelow\fi}

\endinput

% Всё ниже оставлено из исходного файла как справочный образец титульной
% страницы и LaTeX не читает после команды \endinput.
% ============================================================
% Информация о документе и титульная страница
% ============================================================
\date{III семестр, 2026 г.}
\title{Название предмета. Неофициальный конспект}
\author{Лектор: ФИО лектора\\
Конспект: Викторова Дарья, 25.Б03-пу}

\begin{document}

\begin{titlepage}
    \thispagestyle{empty}
    \begin{center}
        \vspace*{0.6cm}

        {\large Санкт-Петербургский государственный университет\par}
        \vspace{0.35em}
        {\normalsize Факультет прикладной математики~--- процессов управления\par}

        \vspace*{5.1cm}

        {\fontsize{30}{35}\selectfont\bfseries Название предмета\par}
        \vspace{0.65em}
        {\fontsize{17}{21}\selectfont\itshape Неофициальный конспект\par}

        \vspace{2.7cm}

        {\normalsize
        \begin{tabular}{@{}rl@{}}
            \textit{Лектор:} & ФИО лектора \\
            \textit{Конспект:} & Викторова Дарья, 25.Б03-пу
        \end{tabular}\par}

        \vfill

        {\large III семестр\par}
        \vspace{0.22em}
        {\large 2026 год\par}
        \vspace*{0.55cm}
    \end{center}
\end{titlepage}

\tableofcontents
\newpage

\setcounter{lection}{0}
\newlection{Число Месяц 2026 г.}

\end{document}



\title{Дифференциальные уравнения. Неофициальный конспект}
\author{Лектор: ФИО лектора\\
Конспект: Викторова Дарья, 25.Б03-пу}
\date{III семестр, 2026 г.}

\begin{document}

\begin{titlepage}
    \thispagestyle{empty}
    \begin{center}
        \vspace*{0.6cm}

        {\large Санкт-Петербургский государственный университет\par}
        \vspace{0.35em}
        {\normalsize Факультет прикладной математики~--- процессов управления\par}

        \vspace*{5.1cm}

        {\fontsize{30}{35}\selectfont\bfseries Дифференциальные уравнения\par}
        \vspace{0.65em}
        {\fontsize{17}{21}\selectfont\itshape Неофициальный конспект\par}

        \vspace{2.7cm}

        {\normalsize
        \begin{tabular}{@{}rl@{}}
            \textit{Лектор:} & Жабко Алексей Петрович \\
            \textit{Конспект:} & Викторова Дарья, 25.Б03-пу
        \end{tabular}\par}

        \vfill

        {\large III семестр\par}
        \vspace{0.22em}
        {\large 2026 год\par}
        \vspace*{0.55cm}
    \end{center}
\end{titlepage}

\tableofcontents
\newpage

\setcounter{lection}{0}

\newlection{5 сентября 2026 г.}

\section*{Литература}
\begin{enumerate}
    \item Жабко А.\ П., Котина Е.\ Д., Чижова О.\ Н. <<Дифференциальные уравнения и устойчивость>>. --- СПб. : Лань, 2015.
    \item Петровский И.\ Г. <<Лекции по теории обыкновенных дифференциальных уравнений>>. --- М. : Физматлит, 2009.
    \item Матвеев Н.\ М. <<Методы интегрирования обыкновенных дифференциальных уравнений>>. --- СПб. : Лань-Пресс, 2003.
    \item Зубов В.\ И. <<Теория колебаний>>. --- М. : Высшая школа, 1979.
    \item Филиппов А.\ Ф. <<Сборник задач по дифференциальным уравнениям>>. --- М. : Интеграл-Пресс, 1998.
   \item Матвеев Н.\ М. <<Сборник задач и упражнений по обыкновенным дифференциальным уравнениям>>. --- М. : Наука, 1976.
\end{enumerate}
\section*{Вспомогательная литература}
\begin{enumerate}
    \item Еругин Н.\ П. <<Книга для чтения по общему курсу дифференциальных уравнений>>. --- Минск : Наука и техника, 1972.
    \item Понтрягин Л.\ С. <<Обыкновенные дифференциальные уравнения>>. --- М. : Наука, 1974.
    \item Зубов В.\ И. <<Колебания и волны>>. --- Л. : Издательство ЛГУ, 1989.
    \item Гюнтер Н.\ М., Кузьмин Р.\ О. <<Сборник задач по высшей математике>>. --- СПб. : Лань, 2022.
\end{enumerate}

\newpage

\chapter{Обыкновенные дифференциальные уравнения первого порядка}

\example{
Функция двух переменных: $f(x,y)=\sin(x+y)$.
\[
\frac{\partial f(x,y)}{\partial x}
= \frac{\partial \sin(x+y)}{\partial x}
= \cos(x+y)
= \frac{\partial \sin(x+y)}{\partial y}.
\]
}

\example{
Функция трёх переменных: $f(x,y,z)=e^z\cdot x + y$.
}

\section{Основные понятия и определения}

\definition{
Дифференциальное уравнение называется \textit{обыкновенным}, если оно содержит производные от искомой функции только по одной переменной.
}

\definition{
\textit{Порядком} обыкновенного дифференциального уравнения называется порядок старшей производной, входящей в это уравнение.
}

\definition{
\textit{Обыкновенным дифференциальным уравнением первого порядка} называется соотношение, связывающее искомую функцию, ее аргумент и первую производную от искомой функции.
\begin{enumerate}
    \item Уравнение в \textit{неявной форме}:
    \[
    f(x,y,y')=0,\qquad(1)
    \]
    где $x$ --- независимая переменная, \\
    $y$ --- искомая функция (зависимая переменная), $y'=\dfrac{\mathrm{d}y}{\mathrm{d}x}$.
    \item Уравнение в явном виде (\textit{нормальной форме}):
    \[
    y'=F(x,y).\qquad (1')
    \]
\end{enumerate}
}

\begin{enumerate}
    \item $f(x,y,z)\in C(\underbrace{D}_{\text{область }D}),\ D\subset \mathbb R^3$.
    \item $F(x,y)\in C(G),\ G\subset \mathbb R^2$.
\end{enumerate}

\definition{
\textit{Решением дифференциального уравнения первого порядка в явной форме} называется функция $y=\varphi(x)\in C^1$ на некотором интервале $(a,b)$, которая удовлетворяет следующим свойствам:
\begin{enumerate}
    \item $(x,\varphi(x))\in G$, $x\in(a,b)$ (график решения лежит в области $G$);
    \item $\forall x\in(a,b)\colon \varphi'(x)\equiv F(x,\varphi(x))$.
\end{enumerate}
}

\example{
Для уравнения $x+y\cdot y'=0$, $D=\mathbb R^3$.
\begin{enumerate}
    \item Попробуем получить $x+y(x)\cdot y'(x)\equiv 0$. \\
    Рассмотрим
    \[
    0\equiv\int_0^x\bigl[\tau+y(\tau)\cdot y'(\tau)\bigr]\,\mathrm{d}\tau.
    \]
    Тогда
    \[
    \int_0^x \tau\,\mathrm{d}\tau+\int_0^x y(\tau)y'(\tau)\,\mathrm{d}\tau
    =\frac{1}{2}x^2+\int_0^{y(x)} y\,\mathrm{d}y
    =\frac{1}{2}x^2+\frac{1}{2}y^2=0.
    \]

    \item Решение в неявном виде:
    \[
    \frac{1}{2}x^2+\frac{1}{2}y^2=C.
    \]
    Отсюда
    \[
    y(x)=\pm\sqrt{2C-x^2}.
    \]
    При $C=\frac{1}{2}$ получим $y(x)=\pm\sqrt{1-x^2}$.

    \item Приведём уравнение к нормальной форме: $y'=-\dfrac{x}{y}$. Тогда
    \[
    \left[
    \begin{array}{ll}
    y>0\colon \varphi(x)=\sqrt{1-x^2},\\
    y<0\colon\varphi(x)=-\sqrt{1-x^2},
    \end{array}
    \right.
    \quad x\in(-1,1).
    \]

    \item Решение может быть записано в \textit{параметрической форме}:
    \[
    \begin{cases}
        x=\cos t,\\
        y=\sin t,
    \end{cases}
    \quad t\in(-\pi,\pi).
    \]
\end{enumerate}
}

% =====================================================================
% Проверка неявно заданной функции как решения
% =====================================================================

Пусть функция $g(x,y)=0$ задана неявно. Хотим проверить, является ли это множество \textit{решением дифференциального уравнения в нормальной форме}:
\[
y'=F(x,y). \qquad (1')
\]

Для этого подставим в уравнение функцию $y=y(x)$, которая неявно задаётся условием $g(x,y(x))\equiv 0$. Продифференцируем это тождество по $x$:
\[
\begin{cases}
    \frac{\partial{g(x,y)}}{\partial x} +\frac{\partial g(x,y)}{\partial y} \cdot y'(x)\equiv 0 \\
    g(x,y)=0
\end{cases}
\]

Из этого уравнения можно выразить $y'=\frac{-\frac{\partial g(x,y)}{\partial x}}{\frac{\partial g(x, y)}{\partial y}}$.

Чтобы $g(x,y)=0$ была решением уравнения $(1')$, необходимо и достаточно, чтобы это выражение совпадало с правой частью $F(x,y)$ во всех точках кривой $g(x,y)=0$. Проверим это условие:

\[
\begin{cases}
y'=\frac{-\frac{\partial g(x,y)}{\partial x}}{\frac{\partial g(x, y)}{\partial y}} = F(x,y) \\
g(x,y)=0
\end{cases}
\]

Если система выполнена, то неявная функция является решением.

\example{Для $y+e^y=x$. \\
Рассмотрим параметрическое представление кривой:
\[
\begin{cases}
    x=\varkappa(t)\\
    y=\chi(t)
\end{cases}
\]
Проверим, является ли она решением дифференциального уравнения $y'=F(x,y)$.
\[
y'=\frac{\mathrm{d}y}{\mathrm{d}x}=\frac{\dot\chi(t)}{\dot\varkappa(t)}.
\]
Тогда условие того, что кривая является решением, запишется как
\[
y'=\frac{\dot\chi(t)}{\dot\varkappa(t)}=F(\varkappa(t),\chi(t)),\quad t\in(\alpha,\beta).
\]
\textit{Параметрическое задание удобно, когда явное выражение $y(x)$ получить затруднительно}.
}

\example{
\begin{enumerate}

    \item Рассмотрим уравнение
    \[
    y'=-\frac{x}{y}.\qquad (*)
    \]
    Положим, $C=1$ (то есть рассмотрим окружность единичного радиуса). \\
    Дифференцирование тождества $x^2+y^2=1$ даёт $2x+2y\cdot y'=0$, откуда $y'=-\frac{x}{y}$. \\
    Следовательно, окружность является \textit{интегральной кривой}. \\
    При $y=0$ правая часть не определена. Следовательно, решение существует на полуокружностях ($y>0$ и $y<0$).


    \item Перевернём уравнение $(*)$ относительно $x$ (они эквивалентны):
    \[
    \frac{\mathrm{d}x}{\mathrm{d}y}=-\frac{y}{x}.\qquad(**)
    \]
    Тогда
    \[
    x=\pm\sqrt{1-y^2}.
    \]
    При $x=0$  правая часть перевёрнутого уравнения не определена, поэтому решение существует лишь на полуокружностях ($x>0$ и $x<0$).

    \begin{center}
    \begin{tikzpicture}[scale=2.5]
        \draw[->] (-1.5,0) -- (1.5,0) node[right] {$x$};
        \draw[->] (0,-1.5) -- (0,1.5) node[above] {$y$};
        \draw[thick] (0,0) circle (1cm);
        \filldraw[white] (1,0) circle (1.2pt);
        \draw (1,0) circle (1.2pt);
        \filldraw[white] (-1,0) circle (1.2pt);
        \draw (-1,0) circle (1.2pt);
        \filldraw[white] (0,1) circle (1.2pt);
        \draw (0,1) circle (1.2pt);
        \filldraw[white] (0,-1) circle (1.2pt);
        \draw (0,-1) circle (1.2pt);
        \node[below right] at (1,0) {$(1,0)$};
        \node[below left] at (-1,0) {$(-1,0)$};
        \node[above right] at (0,1) {$(0,1)$};
        \node[below right] at (0,-1) {$(0,-1)$};
    \end{tikzpicture}
    \end{center}

    \definition{\textit{Интегральная кривая} --- график решения уравнения на плоскости. \\
    В данном случае --- совокупность решений прямого $(*)$ и перевёрнутого $(**)$ уравнений.
    }

    \item Общая формула перевёрнутого уравнения для уравнения $y'=F(x,y)$:
    \[
    \frac{\mathrm{d}x}{\mathrm{d}y}=\frac{1}{F(x,y)}.
    \]

    \item Теперь поставим обратную задачу: пусть задано семейство кривых $y=q(x,c)$. Требуется найти дифференциальное уравнение $y'=F(x,y)$, для которого все кривые этого семейства являются интегральными. \\
    Иными словами, нужно определить функцию $F$ так, чтобы каждая кривая $y=q(x,C)$ при любом $C$ удовлетворяла уравнению $y'=F(x,y)$.

    Для решения заметим: если при некотором фиксированном $C$ кривая $y=q(x,C)$ есть решение, то для всех $x$ должны выполняться тождества
    \[
    \begin{cases}
        y(x)\equiv q(x,C)\,\\
        y'(x)=\frac{\partial q(x,C)}{\partial x}.
    \end{cases}
    \]
    Из соотношения $y=q(x,C)$ выразим $C$ как функцию $x$ и $y$: $C=p(x,y)$. Подставив это выражение в формулу для производной, получим искомое уравнение:
    \[
    y'=F(x,y)=\frac{\partial q(x,C)}{\partial x},\qquad C=p(x,y).
    \]

\end{enumerate}

    \textit{Проиллюстрируем процедуру исключения параметра на примере: \\
    Пусть семейство задано неявно.
    \[
    x^2+y^2-2C=0.
    \]
    Дифференцируем по $x$: $2x+2y\cdot y'=0$, откуда $y'=-\frac{x}{y}$. \\
    В этом примере параметр $C$ в производную не входит, поэтому после дифференцирования он автоматически исключается, и мы сразу получаем дифференциальное уравнение.}

}

\section{Геометрическая и механическая интерпретация}

Рассмотрим уравнение в нормальной форме
\[
y'=f(x,y),\quad f\in C(D).\qquad (1)
\]

\definition{
\textit{Полем направлений} системы $(1)$ называется совокупность отрезков единичной длины, проведённых в каждой точке \((x,y)\in D\) таким образом, что угловой коэффициент (тангенс угла наклона) каждого отрезка к оси \(Ox\) равен значению \(f(x,y)\) в этой точке.
}

\noindent
\textit{Обозначим} через \(\alpha(x,y)\) угол между проведённым отрезком и положительным направлением оси \(Ox\). Тогда
\[
\operatorname{tg}\alpha(x,y)=f(x,y).
\]

\fact{
Направление касательной в каждой точке интегральной кривой совпадает с направлением поля уравнения \(y'=f(x,y)\) в этой точке.
}

\definition{
\textit{Изоклиной} называется кривая, задаваемая уравнением
\[
f(x,y)=C,\qquad C-\text{const}.
\]
В точках изоклины все отрезки поля направлений имеют один и тот же наклон, равный \(\alpha=\operatorname{arctg} C\).
}

\example{
Рассмотрим уравнение $y'=x^2+y^2$.

\[
\textcolor{red}{C=0}\qquad \textcolor{blue}{C=\frac12}\qquad \textcolor{green!50!black}{C=1}\qquad \textcolor{orange}{C=4}.
\]

\begin{center}
\begin{tikzpicture}[scale=1.2]
    % Оси
    \draw[->] (-2.5,0) -- (2.5,0);
    \draw[->] (0,-2.5) -- (0,2.5);

    % Поле направлений для y' = x^2+y^2
    \pgfmathsetmacro{\step}{0.4}
    \pgfmathsetmacro{\len}{0.2}
    \foreach \x in {-2.0,-1.6,...,2.0} {
        \foreach \y in {-2.0,-1.6,...,2.0} {
            \pgfmathsetmacro{\f}{\x*\x+\y*\y}
            \pgfmathsetmacro{\ang}{atan(\f)}
            \pgfmathsetmacro{\dx}{\len*cos(\ang)}
            \pgfmathsetmacro{\dy}{\len*sin(\ang)}
            \draw[->, thin, gray] (\x-\dx/2,\y-\dy/2) -- (\x+\dx/2,\y+\dy/2);
        }
    }

    % Изоклины (окружности) для C = 0, 0.5, 1, 4
    \draw[red, thick] (0,0) circle (0);
    \filldraw[red] (0,0) circle (1.2pt);

    \draw[blue, thick] (0,0) circle ({sqrt(0.5)});
    \draw[green!50!black, thick] (0,0) circle (1);
    \draw[orange, thick] (0,0) circle (2);
\end{tikzpicture}
\end{center}


На рисунке видно, что на каждой окружности все отрезки поля имеют одинаковый угол наклона.\\
Таким образом, изоклины дают наглядное представление о том, как меняется направление поля в зависимости от положения точки.
}

\section{Задача Коши}

Рассмотрим уравнение в нормальной форме
\[
y'=f(x,y),\qquad (x,y)\in D.\qquad (1)
\]

\example{
Рассмотрим модель радиоактивного распада: пусть \(y(t)\) --- масса вещества в момент времени \(t\). Скорость распада пропорциональна массе:
\[
y'=-k y,\qquad k>0.
\]
Общее решение этого уравнения:
\[
y(t)=C e^{-kt}.
\]
Действительно, \(y'(t)=-k C e^{-kt}=-k y(t)\), значит, функция удовлетворяет уравнению при любом \(C\). \\
Если задано начальное условие \(y(0)=y_0\), то \(C=y_0\), и решение принимает вид
\[
y(t)=y_0 e^{-kt}.
\]
В общем случае, при начальном условии \(y(t_0)=y_0\), решение записывается как
\[
y(t)=y_0 e^{-k(t-t_0)}.
\]
}

\definition{
\textit{Задачей Коши} (\textit{начальной задачей}) для уравнения $(1)$ называется задача отыскания решения \(y=\varphi(x)\) этого уравнения, удовлетворяющего дополнительному условию
\[
\varphi(x_0)=y_0,\qquad (x_0,y_0)\in D. \tag{2}
\]
\begin{enumerate}
    \item Числа \(x_0,y_0\) называются \textit{начальными данными},
    \item Условие $(2)$ называется \textit{начальным условием}.
\end{enumerate}
}

\newlection{}

\end{document}

